% from ega3-6-fr.tex
\oldpage[III]{160}
recouvrement $\mathfrak U^{(i)}$ la famille réduite au seul élément $X^{(i)}$,
de sorte qu'avec les notations de (6.6.2), $L_{\bullet\bullet}^{(i)}$ se réduit
à $\Gamma(X^{(i)},\mathscr P_\bullet^{(i)})$ (considéré comme bicomplexe dont
les termes de premier degré $\ne0$ sont nuls), et l'égalité des deux membres
de (6.7.4.2) résulte de (6.4.1.1) et (6.3.1).
\end{remark}

\begin{proposition}[6.7.5]
Soit $u:\mathscr P_\bullet^{(1)}\to\mathscr Q_\bullet^{(1)}$ un
homomorphisme de complexes de $\mathscr O_{X^{(1)}}$-Modules quasi-cohérents,
limités inférieurement, tel que l'homomorphisme
\[
  \mathscr H_\bullet(u):
  \mathscr H_\bullet(\mathscr P_\bullet^{(1)})
  \longrightarrow
  \mathscr H_\bullet(\mathscr Q_\bullet^{(1)})
\]
déduit de $u$ soit un isomorphisme. Alors les homomorphismes
\[
  {}^{(t)}\mathscr E
  (f_1,f_2;\mathscr P_\bullet^{(1)},\mathscr P_\bullet^{(2)})
  \longrightarrow
  {}^{(t)}\mathscr E
  (f_1,f_2;\mathscr Q_\bullet^{(1)},\mathscr P_\bullet^{(2)})
\]
déduits de $u$ sont des isomorphismes pour $t=a$, $t=b$ et $t=c$.
\end{proposition}

\begin{proof}
L'assertion relative à la suite spectrale $(c)$ résulte de ce que cette suite
est birégulière et de ce que l'homomorphisme considéré est un isomorphisme pour
les termes $E_2$ par hypothèse ($0_{\mathrm I}$, 11.1.5). Ceci montre déjà que
\[
  \mathfrak{Tor}_\bullet^S
  (f_1,f_2;\mathscr P_\bullet^{(1)},\mathscr P_\bullet^{(2)})
  \longrightarrow
  \mathfrak{Tor}_\bullet^S
  (f_1,f_2;\mathscr Q_\bullet^{(1)},\mathscr P_\bullet^{(2)})
\]
est un isomorphisme. Appliquant les relations (6.7.4.1) et (6.7.4.2) on voit
d'abord que les homomorphismes
\[
  \mathscr H^{-n}(f_1,\mathscr P_\bullet^{(1)})
  \longrightarrow
  \mathscr H^{-n}(f_1,\mathscr Q_\bullet^{(1)})
  \quad\text{et}\quad
  \mathscr{T}\!or_n^S
  (\mathscr P_\bullet^{(1)},\mathscr P_\bullet^{(2)})
  \longrightarrow
  \mathscr{T}\!or_n^S
  (\mathscr Q_\bullet^{(1)},\mathscr P_\bullet^{(2)})
\]
déduits de $u$ sont des isomorphismes. L'assertion relative aux suites $(a)$
et $(b)$ résulte alors de ce que ces suites sont birégulières (6.7.3) et que
les homomorphismes considérés sont bijectifs pour les termes $E_2$
($0_{\mathrm I}$, 11.1.5).

Notons d'autre part que, si
$u:\mathscr P_\bullet^{(1)}\to\mathscr Q_\bullet^{(1)}$ est un
\emph{homotopisme}, on en déduit des isomorphismes canoniques
\[
  {}^{(t)}\mathscr E
  (f_1,f_2;\mathscr P_\bullet^{(1)},\mathscr P_\bullet^{(2)})
  \longrightarrow
  {}^{(t)}\mathscr E
  (f_1,f_2;\mathscr Q_\bullet^{(1)},\mathscr P_\bullet^{(2)})
\]
pour les six suites spectrales. En effet, si $S$ et les $Y^{(i)}$ sont affines,
on déduit de $u$ un homotopisme de bicomplexes
\[
  C^\bullet(\mathfrak U^{(1)},\mathscr P_\bullet^{(1)})
  \longrightarrow
  C^\bullet(\mathfrak U^{(1)},\mathscr Q_\bullet^{(1)}),
\]
et la proposition résulte de la théorie générale de l'hyperhomologie
($0_{\mathrm I}$, 11.3.2); le passage au cas général se fait par recollement,
en utilisant le fait que, d'un homotopisme de complexes, on déduit un
homotopisme de résolutions projectives de Cartan-Eilenberg de ces complexes
(M, XVII, 1.2).
\end{proof}

\begin{proposition}[6.7.6]
Supposons que le complexe $\mathscr P_\bullet^{(1)}$ ou le complexe
$\mathscr P_\bullet^{(2)}$ soit formé de Modules $S$-plats (les deux complexes
étant limités inférieurement). On a alors un isomorphisme canonique de
bi-$\partial$-foncteurs
\[
  \mathfrak{Tor}_n^S
  (f_1,f_2;\mathscr P_\bullet^{(1)},\mathscr P_\bullet^{(2)})
  \xrightarrow{\sim}
  \mathscr H^{-n}\bigl(f_1\times_S f_2,
    \mathscr P_\bullet^{(1)}\otimes_S\mathscr P_\bullet^{(2)}\bigr).
  \tag{6.7.6.1}\label{III.6.7.6.1.fr}
\]
\end{proposition}

\begin{proof}
Supposons d'abord $S$, $Y^{(1)}$ et $Y^{(2)}$ affines, de sorte qu'on est
dans la situation de (6.6.2), dont nous conservons les notations. Supposons par
exemple que $\mathscr P_\bullet^{(1)}$ soit formé de Modules $S$-plats, et
calculons l'hypertor en utilisant la remarque (6.6.6) : c'est donc l'homologie
de $L_{\bullet\bullet}^{(1)}\otimes_A M_{\bullet\bullet\bullet}^{(2)}$, où
$M_{\bullet\bullet\bullet}^{(2)}$ est une résolution projective de
Cartan-Eilenberg de $L_{\bullet\bullet}^{(2)}$, au sens de
($0_{\mathrm I}$, 11.7.1). D'autre part, les modules
$L_{ij}^{(1)}$ sont \emph{plats} sur $A$ en vertu de l'hypothèse (I, 4.15.1);
on déduit alors de ($0_{\mathrm I}$, 11.7.5) un isomorphisme canonique
\[
  \mathbf{Tor}_\bullet^S
  (\mathfrak U^{(1)},\mathfrak U^{(2)};
   \mathscr P_\bullet^{(1)},\mathscr P_\bullet^{(2)})
  \xrightarrow{\sim}
  \mathscr H_\bullet
  (L_{\bullet\bullet}^{(1)}\otimes_A L_{\bullet\bullet}^{(2)}).
  \tag{6.7.6.2}\label{III.6.7.6.2.fr}
\]
On a d'autre part un homomorphisme naturel de bicomplexes de
$L_{\bullet\bullet}^{(1)}\otimes_A L_{\bullet\bullet}^{(2)}$ dans
$C^\bullet(\mathfrak U,\mathscr Q_\bullet)$, où $\mathfrak U$ est le
recouvrement de $Z=X^{(1)}\times_S X^{(2)}$ par les ouverts affines
